\documentclass[11pt,a4paper]{article}
\usepackage[T1]{fontenc}
\usepackage[utf8]{inputenc}
\usepackage[provide=*,french]{babel}
\usepackage{lmodern,microtype}
\usepackage[a4paper,margin=23mm,headheight=15pt]{geometry}
\usepackage{amsmath,amssymb,booktabs,tabularx,array,enumitem,fancyhdr,xcolor}
\usepackage[colorlinks=true,linkcolor=blue!40!black,urlcolor=blue!40!black]{hyperref}
\usepackage{xurl}
\definecolor{accent}{HTML}{173B5B}
\pagestyle{fancy}\fancyhf{}
\fancyhead[L]{\small\color{accent} P1 -- Note de réorganisation V1.4}
\fancyhead[R]{\small 31 octobre 2026}
\fancyfoot[C]{\thepage}
\renewcommand{\headrulewidth}{0.3pt}
\setlength{\parindent}{0pt}
\setlength{\parskip}{6pt}
\setlength{\emergencystretch}{2em}
\setlist{itemsep=4pt,topsep=4pt,leftmargin=18pt}
\renewcommand{\arraystretch}{1.15}
\newcommand{\E}{\mathbb E}
\newcommand{\Law}{\mathcal L}
\newcommand{\eps}{\varepsilon}
\hypersetup{pdftitle={P1 V1.4 -- Article principal et compagnon technique},pdfauthor={Brice Pouly},pdfsubject={Répartition, conservation des résultats, dépendances et contrôles de la version 1.4}}
\widowpenalty=10000\clubpenalty=10000
\begin{document}
{\Large\bfseries\color{accent} P1 V1.4 : article et compagnon technique}\\[4pt]
{\large Note de réorganisation -- 31 octobre 2026}

\section*{1. La décision éditoriale mise en œuvre}
La V1.4 présente P1 sous la forme de \textbf{deux documents associés} :
un article principal de \textbf{60 pages} et un compagnon technique de
\textbf{93 pages}. La V1.3 réunissait le tout en un manuscrit de 142 pages.
La réduction porte donc sur la longueur du parcours principal ; le volume
total augmente légèrement, à cause des présentations, guides, rappels et
bibliographies nécessaires à la lecture des deux textes.

Les deux documents constituent \textbf{un même travail scientifique}.
Le compagnon contient des preuves nécessaires aux théorèmes de l'article,
ainsi que les prolongements conservés de la V1.3. Il doit accompagner
l'article lors d'une évaluation du travail complet. Cette organisation
ne présente pas le compagnon comme un second article de recherche autonome.

La V1.4 est une réorganisation de la démonstration et de sa présentation.
Elle ne revendique pas de nouvel exposant par rapport à la V1.3.
Les résultats optimaux, les hypothèses d'uniformité et les restrictions
de portée restent ceux de cette version.

\subsection*{Le fil directeur de l'article}
L'article part directement des probabilités marginales exactes
\(p_x=\E J_x\), plutôt que de refaire intégralement l'approximation sur
l'ensemble régulier avant de la raffiner. La rigidité locale est placée
près de cette définition : pour les fenêtres admissibles,
\(p_x\in\{2^{-L},2^{1-L}\}\), à partir d'un seuil suffisamment grand.
Les événements de probabilité nulle sont simplement omis.

Le lecteur suit ensuite les identités affines, l'arrêt héréditaire,
les entrées arithmétiques, le contrôle des triplets, puis la reconstruction
et la borne inférieure. Les calculs longs sont appelés par des énoncés
quantifiés et des renvois précis vers leurs preuves dans le compagnon.

\subsection*{La place du taux optimal}
Avec \(\mu_N=\E Z_N\), \(\alpha_N=\operatorname{Var}(Z_N)-\mu_N\),
\(\Pi_\mu=\operatorname{Pois}(\mu)\) et
\(\Delta=\delta_1-\delta_0\), le théorème 1.1 conserve
\[
 c_{C,\eps}N^{-4/3-\eps}
 \le \left\|\Law(Z_N)-\Pi_{\mu_N}
 -\frac{\alpha_N}{2}\Pi_{\mu_N}*\Delta^{*2}\right\|_1
 \le C_{C,\eps}N^{-4/3+\eps}.
\]
L'optimalité porte sur la puissance de \(N\), avec une marge fixe
arbitraire ; elle ne détermine ni le facteur logarithmique ni la constante
du reste. La preuve reste annoncée sans hypothèse conjecturale.

\clearpage
\section*{2. Une preuve principale suivie sans détour}
Le partage est organisé par dépendances. Le tableau indique ce que le
lecteur trouve dans l'article et ce qu'il consulte dans le compagnon.
Les numéros commençant par T désignent toujours ce dernier.

\small
\begin{tabularx}{\textwidth}{@{}>{\raggedright\arraybackslash}p{.23\textwidth}>{\raggedright\arraybackslash}X>{\raggedright\arraybackslash}X@{}}
\toprule
Étape & Article principal & Compagnon technique\\
\midrule
Modèle et marges
& Modèle, résultats, marges exactes, preuve courte de rigidité locale ; sections 1--2.
& Description des blocs locaux et conséquences supplémentaires ; T19.\\[4pt]
Première interaction
& Mécanisme des canaux rationnels et profil complet des paires ; section 3.
& Estimations arithmétiques, budgets et preuves du profil ; T1--T3.\\[4pt]
Algèbre et arrêt
& Identités centrées, composantes directes, arrêt exact et exemple de suffixe dépendant ; sections 4--5.
& Classe retenue de la construction sur l'ensemble régulier ; T12.\\[4pt]
Sources et uniformité
& Entrées quantitatives, moments fixés indépendants, certificats uniformes ; sections 6--10.
& Géométrie des sources, moments de sélections courtes, itérations et certificats ; T4--T6 et T13.\\[4pt]
Triplets et rigidité
& Identités signées, rang des paires, courbe commune, grille, exposant \(41/30\) ; sections 11--14.
& Partitions, confinement, construction biquadratique, résultants, noyaux résiduels ; T7--T10.\\[4pt]
Reconstruction
& Réduction des ordres fixés, queue uniforme, arrêt et borne inférieure quartique ; sections 15--16.
& Deux preuves longues de dénombrement conditionnel ; T11.\\[4pt]
Conséquences
& Poids polynomiaux, canaux explicites, Hermite et énoncé positionnel ; sections 17--20.
& Normes, preuve positionnelle, suppression, queues et limites de méthode ; T18--T23.\\
\bottomrule
\end{tabularx}
\normalsize

\subsection*{Ce qui reste visible dans l'article}
La partie arithmétique conserve les idées qui expliquent les exposants :
extraction de relations courtes, existence d'une courbe commune dont les
bornes ne dépendent pas du nombre de vecteurs de base, passage à la grille
des décalages et estimation de points entiers uniforme en degré.
Le théorème 13.6 conserve le seuil suffisant
\((2/3)\log 2<1/2\), sans le déclarer optimal.

La géométrie du triplet \(10:12:15\), l'inégalité de recouvrement
\(19/30\), sa compatibilité affine et la preuve du cumulant d'ordre trois
restent également dans l'article. Le gain
\(41/30-4/3=1/30\) suffit au majorant de la reconstruction.
La borne inférieure du troisième cumulant est une conclusion distincte ;
elle n'est pas utilisée pour prouver ce majorant.

\clearpage
\section*{3. Correspondance avec la V1.3}
La table complète de l'archive permet de retrouver chaque ancien repère.
Le tableau ci-dessous fournit le déplacement des grands blocs.

\small
\begin{tabularx}{\textwidth}{@{}>{\raggedright\arraybackslash}p{.20\textwidth}>{\raggedright\arraybackslash}p{.28\textwidth}>{\raggedright\arraybackslash}X@{}}
\toprule
V1.3 & Article V1.4 & Compagnon V1.4\\
\midrule
Sections 1--2 & Sections 1 et 3 & Preuves de profil en T3, entrées en T1--T2.\\[4pt]
Sections 3--9 & Outils réutilisés en 4--8 & Preuves à sources fixées en T4 ; construction complète sur l'ensemble régulier en T12--T17.\\[4pt]
Section 10 & Section 2 & La rigidité locale provient aussi de l'ancien lemme F.2.\\[4pt]
Sections 11--14 & Sections 9--12 & Preuves techniques en T5--T8.\\[4pt]
Section 15 & Section 13 & Extraction, construction et résultants en T9.\\[4pt]
Section 16 & Section 14 & Triplets au-dessus des canaux alignés en T10.\\[4pt]
Section 17 & Section 15 & Dénombrements denses et petits noyaux en T11.\\[4pt]
Sections 18--20 & Sections 16--18 & Les preuves et conclusions centrales restent dans l'article.\\[4pt]
Section 21 & Section 19 & Preuve des normes de différences en T18.\\[4pt]
Annexes A--B & Énoncés utiles rappelés aux interfaces & T1--T2.\\[4pt]
Annexes C--E & Renvois de portée & T20--T22.\\[4pt]
Annexe F & Rigidité locale en 2 ; énoncé positionnel en 20 & Blocs, queues, preuve positionnelle, Palm, suppression et obstructions en T19.\\[4pt]
Annexe G & Limite de la méthode rappelée dans le texte & Contre-exemple complet en T23.\\
\bottomrule
\end{tabularx}
\normalsize

\subsection*{Conservation vérifiée}
Les \textbf{505 repères} de la V1.3 sont conservés chacun une fois.
Les deux textes comportent toujours \textbf{167 énoncés} de type théorème,
proposition, lemme ou corollaire : 73 dans l'article et 94 dans le compagnon.
Les 166 environnements de démonstration sont également conservés dans
l'ensemble. Quarante nouveaux repères servent à la navigation et aux
interfaces de preuve.

Ce contrôle inclut les identifiants et les déplacements ; il ne constitue
pas à lui seul une validation mathématique. Les adaptations des identités
affines aux marges exactes, les modifications de notation et les ajustements
de preuves correspondants sont documentés dans les comptes rendus de sources.

\subsection*{Pourquoi l'ancienne construction est conservée}
Le taux \(N^{-4/5+\eps}\) de la reconstruction sur l'ensemble régulier
reste utile pour le champ qui garde les positions et pour la suppression
des doublets. Le déplacer dans le compagnon évite de faire parcourir
deux reconstructions complètes au lecteur du résultat scalaire optimal,
tout en préservant les outils nécessaires aux prolongements.

\clearpage
\section*{4. Résultats conservés et limites explicites}
Les précisions ci-dessous concernent des références différentes. Leur
distinction est maintenue dans les énoncés et dans l'introduction.

\small
\begin{tabularx}{\textwidth}{@{}>{\raggedright\arraybackslash}p{.28\textwidth}>{\raggedright\arraybackslash}p{.23\textwidth}>{\raggedright\arraybackslash}X@{}}
\toprule
Objet & Conclusion & Emplacement dans l'article\\
\midrule
Poisson de moyenne nominale ou exacte
& Équivalent d'ordre \(LN^{-2/3}\), coefficient explicite
& Corollaire 1.2. Dépendance en \(N2^{-L}\) et \(L\bmod 3\).\\[4pt]
Correction quadratique avec \(\mu_N,\alpha_N\) exacts
& Puissance optimale \(4/3\)
& Théorème 1.1 ; reconstruction 15.14 ; borne inférieure 16.2 ; assemblage en 17.\\[4pt]
Troisième cumulant factoriel
& Puissance optimale \(41/30\)
& Théorème 14.1. Famille compatible \(10:12:15\), sans constante asymptotique générale annoncée.\\[4pt]
Canaux jusqu'à une hauteur \(h\) fixée
& Puissance optimale \(h/(h+1)\) et équivalent du premier niveau omis
& Théorème 18.1. En particulier \(3/4,4/5,5/6\) pour \(h=3,4,5\).\\[4pt]
Poids polynomiaux et moments fixés
& Reste \(O(N^{-4/3+\eps})\)
& Corollaire 17.1 ; ordre de moment fixé avant \(N\).\\[4pt]
Absence de début retenu
& \(e^{-\mu_N}(1+\alpha_N/2)+O(N^{-4/3+\eps})\)
& Corollaire 17.2. Les fenêtres peuvent dépasser l'extrémité droite de l'intervalle.\\[4pt]
Référence de Hermite exacte
& Majorant \(O(N^{-4/3+\eps})\)
& Corollaire 19.1. La borne inférieure de la référence quadratique ne lui est pas attribuée.\\[4pt]
Champ positionnel non signé
& Comparaison pondérée à puissance \(4/5\)
& Lemme 20.1 ; preuve en T19. Positions conservées, signes et longueurs excédentaires non conservés.\\
\bottomrule
\end{tabularx}
\normalsize

\subsection*{Précautions qui structurent la preuve}
La moyenne nominale \(\lambda_N=N2^{-L}\) ne remplace pas gratuitement
\(\mu_N\) à la précision \(4/3\) : la proposition 17.3 explicite le
biais quantifié. L'estimation de conversion de puissance \(5/3\)
concerne la référence de Bernoulli \emph{avec tous les termes de covariance}.

Les ordres fixés jusqu'à 256 sont séparés d'une queue uniforme sur les
ordres plus grands. Le nombre 256 est un seuil suffisant, sans prétention
d'optimalité. Les théorèmes de hauteur finie ne sont pas appliqués à
\(h=h(N)\). Enfin, les queues d'un compte modifié ne deviennent pas
automatiquement celles du compte original : les obstructions de T19 et
T23 restent disponibles avec leurs démonstrations.

\clearpage
\section*{5. Contrôles et contenu de la livraison}
La vérification de cette édition porte sur la fidélité de la scission,
la cohérence des interfaces et la qualité des documents. Elle ne remplace
pas une expertise humaine indépendante de l'ensemble des preuves.

\subsection*{Contrôles réalisés}
\begin{itemize}
\item Conservation des 505 repères d'origine, absence de doublon et de
renvoi non résolu ; conservation des 167 énoncés et des 166 environnements
de preuve.
\item Vérification des cibles des \textbf{209 liens entre documents}
et des 608 liens internes aux deux PDF. Les numéros T restent utilisables
lorsque le lecteur de PDF ne suit pas les liens entre fichiers.
\item Relecture des dépendances : moments fixés disponibles avant
l'interpolation, activité des triplets estimée avant reconstruction,
preuve scalaire optimale indépendante de l'ancienne reconstruction globale.
\item Contrôle visuel de toutes les pages, complété par l'examen en
pleine page des tableaux, équations denses et interfaces. Aucune boîte
débordante ni référence manquante dans les compilations finales.
\item Réexécution des vérifications finies pertinentes pour les marges
hétérogènes et l'arrêt : huit systèmes affines, 112 vérifications exactes,
et les identités rationnelles du développement arrêté. Les trois autres
diagnostics arithmétiques de la V1.3 et leurs résultats sont conservés
sans présenter leur exécution antérieure comme une nouvelle certification.
\end{itemize}

La relecture a notamment clarifié la distinction entre le rang brut
\(\rho\), qui compte les défauts locaux, et le rang relatif effectif
\(\rho_{\mathrm{rel}}\). L'intensité incidente du champ positionnel
est notée \(\iota_x\), pour ne pas réutiliser \(d_x=\dim S_x\).
Quelques commandes typographiques héritées ont également été réparées.

\subsection*{Les fichiers livrés}
\begin{itemize}
\item \textbf{Article principal en anglais :} 60 pages, sections 1--20,
résultats, mécanisme arithmétique, reconstruction et optimalité.
\item \textbf{Compagnon technique en anglais :} 93 pages, sections T1--T23,
preuves détaillées et prolongements.
\item \textbf{La présente note en français :} décision éditoriale,
correspondance et portée.
\item \textbf{Archive des sources :} sources des trois PDF, commandes
de compilation, diagnostics, comptes rendus de contrôle et deux tableaux
CSV de correspondance, l'un pour les 505 repères, l'autre pour les
167 énoncés.
\end{itemize}

\subsection*{Reproduction et circulation}
Le fichier \texttt{README.txt} indique les dépendances et les commandes.
Depuis la racine des sources, \texttt{python3 build.py} compile les deux
documents jusqu'à stabilisation des renvois. Le script de la note est
\texttt{editorial/build\_note.sh}. Le contrôle des renvois est relançable
avec \texttt{python3 editorial/check\_references.py}.

Les deux PDF scientifiques portent la même version et la même date.
Ils doivent être diffusés ensemble, en conservant leurs noms de fichiers
si l'on veut utiliser les liens croisés. La V1.3 reste un état cumulatif
de référence ; la V1.4 fournit la présentation organisée pour la lecture
et l'évaluation.
\end{document}
